% Generated draft: requires mathematical and source review.
\documentclass[11pt]{article}
\usepackage[margin=25mm]{geometry}
\usepackage{fontspec}
\setmainfont{Latin Modern Roman}
\usepackage{amsmath,amssymb,mathtools,unicode-math}
\setmathfont{Latin Modern Math}
\usepackage{enumitem,needspace,xurl,hyperref}
\setlength{\parindent}{0pt}
\setlength{\parskip}{5pt}
\setlength{\emergencystretch}{3em}
\newcommand{\N}{\mathbb{N}}
\newcommand{\Z}{\mathbb{Z}}
\newcommand{\Q}{\mathbb{Q}}
\newcommand{\R}{\mathbb{R}}
\newcommand{\C}{\mathbb{C}}
\newcommand{\F}{\mathbb{F}}
\newcommand{\A}{\mathbb{A}}
\newcommand{\PP}{\mathbb{P}}
\newcommand{\E}{\mathbb{E}}
\newcommand{\eps}{\varepsilon}
\newcommand{\dd}{\,\mathrm{d}}
\begin{document}
\section*{UnsolvedMath ID: 20001971}
\textbf{Source problem code:} catalogue entry
\subsection{Definitions and mathematical statement}
R. Goldin's Question 7.3 asks, for a previously specified \(G\)-action on \(V=\C^n\) and a previously defined \(X\), whether
\[
 H_G^*(V)\longrightarrow H^*(X/G)=H_G^*(X)
\]
is surjective, and asks for its kernel. The preceding definitions are absent. The aid reports an internal ambiguity between a locally-free locus before quotienting and its complexified quotient; the displayed equivariant restriction supports the former interpretation, which is adopted only explicitly and conditionally here:
\[
 X=V_{\mathrm{lf}}=\{v\in V:G_v\text{ is finite}\}.
\]
For a compact group the linear contraction gives \(H_G^*(V;R)\cong H^*(BG;R)\). The Borel target is \(H_G^*(X;R)=H^*(EG\times_GX;R)\). Its identification with ordinary cohomology of the coarse quotient is coefficient sensitive for finite stabilizers and cannot simply be assumed integrally.

The aid's specific result concerns every diagonal circle representation
\[
 V=\C^{n_0}\oplus\bigoplus_{i=1}^m\C_{w_i},\qquad
 n_0\ge0,\quad w_i\in\Z\setminus\{0\},
\]
where \(z\in S^1\) acts on \(\C_{w_i}\) by multiplication by \(z^{w_i}\) and acts trivially on \(\C^{n_0}\). If \(m=0\), the locally-free locus is empty. If \(m\ge1\), it is
\[
 X=\C^{n_0}\times(\C^m\setminus\{0\}).
\]
Writing \(H^*(BS^1;\Z)=\Z[u]\), \(\deg u=2\), and using the complex orientation, the moving equivariant Euler class is
\[
 e=\Bigl(\prod_{i=1}^m w_i\Bigr)u^m.
\]
The reported restriction is surjective over both \(\Z\) and \(\Q\), with
\[
 H_{S^1}^*(X;\Z)\cong\Z[u]/(e),\qquad
 \ker\bigl(\Z[u]\to H_{S^1}^*(X;\Z)\bigr)=(e),
\]
and
\[
 H_{S^1}^*(X;\Q)\cong\Q[u]/(u^m),\qquad
 \ker\bigl(\Q[u]\to H_{S^1}^*(X;\Q)\bigr)=(u^m).
\]
No effectiveness or primitivity of the weights is assumed. If the nonzero moving coordinates have index set \(I\), the stabilizer has order \(\gcd_{i\in I}|w_i|\); these are finite even when nontrivial. Rationally the usual locally-free comparison gives the coarse-quotient cohomology, but the displayed integral formula is for the Borel target. Already a single moving weight with \(|w_1|>1\) gives finite-isotropy torsion in the Borel target absent from the contractible radial coarse quotient. For \(m=0\), restriction to the empty set has zero target and kernel the entire coefficient polynomial ring, rather than the nonempty-locus formula.

M. Libine's Comment 7.4 recommends beginning with irreducible representations and warns that \(X\) may be empty without further action hypotheses. The broad arbitrary-group question is not settled by this circle result; the aid says it remained unresolved in the literature it checked. The final Higgs-bundles heading starts the next section.
\subsection{Short English statement}
The aid answers the restriction-and-kernel question for the locally-free locus of any diagonal circle representation. The kernel is the moving Euler ideal, with finite-isotropy torsion retained integrally. The original definition of the target is missing and ambiguous, and rational coarse-quotient cohomology must be distinguished from integral Borel cohomology.
\subsection{Sources}
\begin{itemize}
\item Local UnsolvedMath registry, permanent ID 20001971, exact source code catalogue entry; canonical \texttt{aim-geometry-notes.json}, \texttt{notes[308]}, tag \texttt{question}; AIM Workshop Problem Lists, Geometry. Created and updated 14 August 2026.
\item AIM 2004 workshop list \emph{Moment maps and surjectivity in various geometries}, Question 7.3 (R. Goldin), Comment 7.4 (M. Libine); \url{https://aimath.org/WWN/momentmaps/momentmaps.pdf}. The preceding definitions are not supplied and the PDF was not accessed. No separate authoritative theorem citation is supplied for the Euler-ideal result.
\item Supplied attempt 1: \texttt{partial\_result}, problem status at run \texttt{unknown}, rigor label \texttt{proved}, novelty confidence \texttt{low}; classification \texttt{PARTIAL-PROGRESS}, registry status \texttt{partially\_solved}. Candidate: circle locally-free-locus Euler-kernel theorem, including zero weights, nonprimitive weights, and coarse-quotient distinctions. The arbitrary-group literature assessment is dated supplied evidence, not a new search result.
\item Ulam AI / UnsolvedMath Contributors; CC BY 4.0. Machine-generated research aid, not peer reviewed; independently verify before citation.
\end{itemize}
\end{document}
